\chapter{Methodology}\label{methodology}

\section{Research Design and Overall Methodological Framework}\label{research-design-and-overall-methodological-framework}

The research adopts a quantitative, computational, comparative, empirical, and retrospective design to investigate Bayesian survival modelling for explainable time-to-event prediction in healthcare. Three existing healthcare survival datasets, i.e. GBSG2, METABRIC and WHAS500, are utilized in the empirical investigation and offer patient-level data on survival time, event status and predictors. The retrospective design is suitable since the study does not involve a prospective clinical study or gathering of new patient information but analyzes observations previously made.

The methodology was developed to provide accurate, Bayesian and interpretable results. Firstly, to compare the performance of different survival modelling techniques. Secondly, to examine how specifying different priors affects the posterior estimation of parameters. Finally, to provide uncertainty sensitive interpretations of covariate effects via the posterior distribution. This is in line with recent papers in survival modelling which have shown that model performance can depend on the characteristics of the dataset, outcome structure and assumptions made in different modelling approaches \cite{ref33,ref44}.

\begin{figure}[p]
\centering
\begingroup
\singlespacing
\definecolor{raviblue}{HTML}{CDDEFC}
\definecolor{ravigreen}{HTML}{B4EED1}
\definecolor{raviyellow}{HTML}{FFF5B4}
\definecolor{ravipurple}{HTML}{E0D9FC}
\definecolor{ravipink}{HTML}{FFD0D0}
\definecolor{ravistart}{HTML}{6298EB}
\resizebox{0.96\linewidth}{!}{%
\begin{tikzpicture}[
  x=1cm,y=1cm,
  base/.style={draw=black!80,rounded corners=2pt,line width=0.55pt,
    align=center,font=\sffamily\fontsize{7.0}{7.5}\selectfont,inner sep=2pt,
    text width=4.5cm,minimum height=0.44cm,execute at begin node={\hyphenpenalty=10000}},
  blue/.style={base,fill=raviblue},
  green/.style={base,fill=ravigreen},
  yellow/.style={base,fill=raviyellow},
  purple/.style={base,fill=ravipurple},
  pink/.style={base,fill=ravipink},
  smallgreen/.style={green,text width=3.18cm,minimum height=0.57cm},
  smallpurple/.style={purple,text width=3.18cm,minimum height=0.57cm},
  arrow/.style={-{Stealth[length=1.7mm]},draw=black!75,line width=0.6pt},
  line/.style={draw=black!75,line width=0.6pt},
  edge label/.style={font=\fontsize{7.2}{8}\selectfont,fill=white,inner sep=1pt}
]
% The arrangement follows Ravi's supplied conceptual diagram.
\node[base,ellipse,fill=ravistart,text width=2.7cm,minimum height=0.6cm]
  (start) at (2.4,0) {Start: Research Methodology};
\node[blue] (design) at (2.4,-0.79)
  {Research Design: Quantitative, Computational and Comparative Study};
\node[blue] (focus) at (2.4,-1.56)
  {Bayesian Survival Modelling:\\Prior Specification and Explainability};
\node[green] (cohorts) at (2.4,-2.33)
  {Survival Cohorts: GBSG2, METABRIC, WHAS500};
\node[green] (data) at (2.4,-3.04)
  {Data Acquisition and Cohort Reconstruction};
\node[green] (checks) at (2.4,-3.78)
  {Cohort Checks: Outcomes and Events\\Survival Times and Predictors};
\foreach \a/\b in {start/design,design/focus,focus/cohorts,cohorts/data,data/checks}
  \draw[arrow] (\a.south) -- (\b.north);

\node[green] (preprocessing) at (7.45,-0.69)
  {Training-Fold Preprocessing\\Applied to Held-Out Data};
\node[yellow] (protocol) at (7.45,-1.56)
  {Fixed Models, Seeds and Design\\Prespecified Evaluation Rules};
\node[yellow] (validation) at (7.45,-2.33)
  {Outer Validation: $5\times5$ Folds\\Three Inner Folds: RSF and DeepSurv};
\node[blue] (development) at (7.45,-3.04) {Parallel Model Development};
\draw[arrow] (checks.east) -- (4.95,-3.78) -- (4.95,-0.79) -- (preprocessing.west);
\foreach \a/\b in {preprocessing/protocol,protocol/validation,validation/development}
  \draw[arrow] (\a.south) -- (\b.north);

\node[pink] (classical) at (2.4,-4.73) {Classical Survival Baseline};
\node[purple] (machine) at (7.45,-4.73) {Machine Learning Baselines};
\node[blue] (bayesian) at (12.5,-4.73) {Bayesian Survival Models};
\draw[line] (development.south) -- (7.45,-4.26);
\draw[arrow] (7.45,-4.26) -| (classical.north);
\draw[arrow] (7.45,-4.26) -- (machine.north);
\draw[arrow] (7.45,-4.26) -| (bayesian.north);

\node[pink] (cox) at (2.4,-5.45) {Cox Proportional Hazards Model};
\node[purple] (rsf) at (7.45,-5.45) {Random Survival Forest};
\node[purple] (deep) at (7.45,-6.18) {DeepSurv};
\node[blue] (piecewise) at (12.5,-5.45)
  {Fully Bayesian Piecewise-Exponential\\Proportional Hazards Model};
\node[blue] (likelihood) at (12.5,-6.25)
  {Expanded Poisson Full Likelihood\\Six Baseline Hazard Intervals};
\node[green] (baselinefit) at (2.4,-6.95) {Model Fitting and Prediction};
\draw[arrow] (classical.south) -- (cox.north);
\draw[arrow] (cox.south) -- (baselinefit.north);
\draw[line] (machine.south) -- (7.45,-5.08) -- (4.97,-5.08) -- (4.97,-6.18);
\draw[arrow] (4.97,-5.45) -- (rsf.west);
\draw[arrow] (4.97,-6.18) -- (deep.west);
\draw[line] (rsf.east) -- (9.93,-5.45) -- (9.93,-6.95);
\draw[line] (deep.east) -- (9.93,-6.18);
\draw[arrow] (9.93,-6.95) -- (baselinefit.east);
\draw[arrow] (bayesian.south) -- (piecewise.north);
\draw[arrow] (piecewise.south) -- (likelihood.north);

\node[blue] (priors) at (7.45,-7.71) {Alternative Coefficient Priors};
\draw[arrow] (likelihood.south) -- (12.5,-7.22) -| (priors.north);
\node[blue] (normal) at (2.4,-8.47) {Normal Prior};
\node[blue] (horseshoe) at (7.45,-8.47) {Regularised Horseshoe Prior};
\node[blue] (spike) at (12.5,-8.47) {Continuous Spike-and-Slab Prior};
\draw[line] (priors.south) -- (7.45,-8.1);
\draw[arrow] (7.45,-8.1) -| (normal.north);
\draw[arrow] (7.45,-8.1) -- (horseshoe.north);
\draw[arrow] (7.45,-8.1) -| (spike.north);

\node[yellow] (posterior) at (7.45,-9.28) {Bayesian Posterior Inference};
\draw[line] (normal.south) -- (2.4,-8.86) -- (12.5,-8.86) -- (spike.south);
\draw[line] (horseshoe.south) -- (7.45,-8.86);
\draw[arrow] (7.45,-8.86) -- (posterior.north);
\node[yellow] (eligibility) at (7.45,-10.03) {MCMC Diagnostics and Model Eligibility};
\node[yellow] (diagnostics) at (7.45,-10.72)
  {Sampling Diagnostics: R-hat and ESS\\Divergences, Tree Depth and BFMI};
\draw[arrow] (posterior.south) -- (eligibility.north);
\draw[arrow] (eligibility.south) -- (diagnostics.north);
\node[base,diamond,aspect=3.8,fill=ravigreen,text width=2.65cm,
  inner sep=1pt,minimum height=0.72cm]
  (gate) at (7.45,-11.72) {Does the Model Meet\\Diagnostic Criteria?};
\draw[arrow] (diagnostics.south) -- (gate.north);
\node[yellow] (retain) at (2.4,-11.72) {Retain Eligible Bayesian Results};
\node[yellow] (record) at (12.5,-11.72)
  {Record Diagnostic Outcome; Exclude Non-Eligible Fit from Inference and Prediction};
\draw[arrow] (gate.west) -- node[edge label,above]{Yes} (retain.east);
\draw[arrow] (gate.east) -- node[edge label,above]{No} (record.west);

\node[yellow] (ppc) at (7.45,-12.68)
  {Posterior Predictive Interpretation\\and Prior-Sensitivity Assessment};
\draw[arrow] (retain.south) -- (2.4,-12.68) -- (ppc.west);
\node[green] (evaluation) at (7.45,-13.52) {Held-Out Predictive Evaluation};
\draw[arrow] (ppc.south) -- (evaluation.north);
\draw[arrow] (baselinefit.west) -- (-0.18,-6.95) -- (-0.18,-13.52) -- (evaluation.west);

\node[smallgreen] (concordance) at (1.65,-14.58)
  {Discrimination: Concordance Index};
\node[smallgreen] (brier) at (5.5,-14.58)
  {Prediction Error: Brier Score and Integrated Brier Score};
\node[smallgreen] (auc) at (9.35,-14.58)
  {Time-Dependent Discrimination:\\Time-Dependent AUC};
\node[smallgreen] (calibration) at (13.2,-14.58)
  {Calibration: Survival Calibration Assessment};
\draw[line] (evaluation.south) -- (7.45,-14.03);
\foreach \n in {concordance,brier,auc,calibration}
  \draw[arrow] (7.45,-14.03) -| (\n.north);
\node[green] (combined) at (7.45,-15.63) {Combined Model Evaluation};
\foreach \n in {concordance,brier,auc,calibration}
  \draw[line] (\n.south) -- (\n.south |- 0,-15.12) -- (7.45,-15.12);
\draw[arrow] (7.45,-15.12) -- (combined.north);
\node[purple] (explain) at (7.45,-16.43) {Bayesian Explainability};
\draw[arrow] (combined.south) -- (explain.north);

\node[smallpurple] (coefficients) at (1.65,-17.43) {Posterior Coefficient Distributions};
\node[smallpurple] (intervals) at (5.5,-17.43) {Credible Intervals and Hazard Ratios};
\node[smallpurple] (direction) at (9.35,-17.43) {Posterior Directional Probabilities};
\node[smallpurple] (shap) at (13.2,-17.43) {Exact Linear SHAP Rankings and Stability};
\draw[line] (explain.south) -- (7.45,-16.91);
\foreach \n in {coefficients,intervals,direction,shap}
  \draw[arrow] (7.45,-16.91) -| (\n.north);
\node[green] (synthesis) at (7.45,-18.47) {Comparative Results Synthesis};
\foreach \n in {coefficients,intervals,direction,shap}
  \draw[line] (\n.south) -- (\n.south |- 0,-17.98) -- (7.45,-17.98);
\draw[arrow] (7.45,-17.98) -- (synthesis.north);
\node[green] (interpretation) at (7.45,-19.27)
  {Interpretation against Research Questions and Objectives};
\node[green] (coverage) at (7.45,-20.07)
  {Evaluation Coverage and Generalisability};
\node[base,fill=ravistart] (output) at (7.45,-20.9)
  {Final Output: Evidence-Based Conclusions and Recommendations};
\foreach \a/\b in {synthesis/interpretation,interpretation/coverage,coverage/output}
  \draw[arrow] (\a.south) -- (\b.north);
\end{tikzpicture}%
}
\endgroup
\caption{Conceptual overview of the study design, analysis stages, and links between modelling and interpretation. Predictive checks are computed for each sampled attempt; posterior interpretation and held-out evaluation use accepted Bayesian fits.}
\label{_Toc245464042}
\end{figure}


There are two main areas to the modelling. The first area consists of three different survival models to compare; Cox proportional hazards regression, Random Survival Forest and DeepSurv. The CoxPH model is a traditional statistical model, whereas the RSF and DeepSurv models are non-linear, machine learning and deep learning models. In similar healthcare studies it has been shown that the predictive capabilities of these different models can vary greatly depending on the clinical dataset and analysis scenario \cite{ref33,ref37}.

The second area consists of the fully Bayesian piecewise exponential proportional hazards model. There are three different prior distributions for the coefficients; Normal, Horseshoe and Spike-and-Slab. The piecewise exponential model allows for the baseline hazard to vary over time (i.e. each time interval has a different value). The piecewise exponential model was chosen due to its flexibility in how the baseline hazard is represented \cite{ref22,ref38}.

The use of a fully Bayesian model allows the analysis to describe uncertainty alongside point estimates. Bayesian inference is a type of inference that combines prior distributions with the likelihood to obtain posterior distributions. These posterior distributions can then be used to measure the uncertainty of the parameters. This means that the model will produce posterior distributions for the coefficients, credible intervals, hazard ratio transformations and posteriors for the direction. Bayesian models can also be used to directly compare previous assumptions and how they affect inference \cite{ref40}.

The three different Bayesian families are used for a variety of different methods. The Normal distribution is a standard choice for regularisation. The Horseshoe distribution is a global-local shrinkage method. This has the capability to robustly shrink down the weak signals but keep the relatively robust effects unconstrained. This has been analysed in both regression and survival scenarios within Bayesian models where there can be many predictors with a small signal. The continuous Spike-and-Slab specification allows coefficients to receive different amounts of regularisation. In recent survival studies the spike-and-slab Cox models have been looked at specifically and shown to be used to detect sparse covariates \cite{carvalho2010,piironen2017,ref42,ref43}.

A variety of different metrics have been used to measure the survival specific predictions. The metrics chosen were; concordance index (C-index), Integrated Brier Score (IBS), time dependent AUC and calibration intercepts and slopes. This is an important factor as they provide a variety of different characteristics to measure for survival predictions. Other recent analyses have used the C-index and Brier based metrics to differentiate between the predictive discrimination and prediction error \cite{ref26,ref27}.

The research design is comparative and associational rather than causal. The aim is not to determine that a given clinical variable leads to a given survival outcome. Rather, the study investigates predictive relationships, posterior uncertainty, model behaviour, and the effects of prior specification using retrospective survival data. This difference is significant as the results of the analysis are not to be viewed as the causal clinical influences but as the statistical relationships and predictive values.

In general, the study design combines traditional statistical survival analysis, machine-learning survival prediction, Bayesian probability modelling, posterior diagnostics, posterior predictive evaluation, predictive evaluation and Bayesian explainability into a single empirical model. Explicit model-quality criteria link the computational diagnostics to the interpretation of posterior results.

\section{Research workflow}\label{research-workflow}

The research pipeline involves creating the dataset, fitting the model, evaluating the model, Bayesian interpretability, prior sensitivity and verifying computational integrity. The workflow was created to keep the model building and model acceptance separate as fitting a Bayesian model does not mean that the posterior samples are reliable enough to be interpreted. A structured Bayesian workflow should contain model building, computation, model checking, validation and model improvement \cite{ref53}.

The first stage involved preparing and reconstructing the GBSG2, METABRIC, and WHAS500 survival datasets. The variables used in the modelling were determined and converted to a standard representation of survival-analysis in the form of explanatory variables, observed survival time and event status. This established a consistent analytical framework for comparing the modelling approaches while preserving the characteristics of each cohort.

The second phase defined the validation architecture. Five repeats of five-fold cross-validation were completed for each cohort, resulting in 75 outer folds across the three cohorts. The six model arms were evaluated within this structure, giving 450 model--fold outcomes. Predictive comparisons use paired fold-level results, with uncertainty adjusted for the overlap between training sets.

The third stage of the process broke down the modelling into a baseline and Bayesian approach. CoxPH, RSF and DeepSurv were applied to create a reference survival model. DeepSurv is a neural network version of the Cox proportional hazards predictor. Studies have shown that the relative performance can vary between different datasets. Hence, multiple comparison families should be used \cite{ref4,ref33,ref37}.

The Bayesian stream was based on piecewise-exponential proportional-hazards formulation and the priors were the Normal, Horseshoe and Spike-and-Slab priors. NUTS-based MCMC with 2,000 warm-up iterations, four chains, and 2,000 iterations of the posterior draw was used to perform Bayesian inference. The core sampling configuration was fixed according to the implemented experimental design, while the target acceptance probability varied by prior specification.

The fourth stage involved checking the MCMC diagnostics. This included checking for divergences, rank normalised r-hat, bulk and tail effective sample sizes, BFMI, tree depths and posterior finiteness. Multiple diagnostics were used to ensure that there is no single convergence statistic that can be used to show that the MCMC is behaving properly. The rank normalised diagnostics make the diagnostic more sensitive to sampling issues that may not be picked up by other summaries \cite{ref39}.

The diagnostic criteria were applied to each Bayesian model--fold fit. A fit that did not pass the initial assessment could undergo two specified refits, with increased warm-up and then a non-centred baseline parameterisation. The likelihood, coefficient priors and acceptance criteria remained fixed. Only accepted fits were passed onto posterior interpretation and prediction.

Once the Bayesian models had been kept, posterior predictive checks were performed. Posterior predictive checks check whether data generated from the fitted posterior distribution can replicate the important features of the observed data. Posterior predictive checks can be an important part of Bayesian model checking and should be considered with parameter estimation and computational diagnostics \cite{ref22,ref53}.

The next step was held-out predictive assessment using concordance, Integrated Brier Score, time-dependent AUC and calibration intercepts and slopes. All six model arms were evaluated across the three cohorts. Fold coverage is reported for each model so that the predictive summaries and paired contrasts can be interpreted with their diagnostic acceptance status.

The next step was on Bayesian explainability. In the case of retained Bayesian fits, the means of the posterior coefficients, 95 percent credible intervals, estimates and credible interval estimates of the hazard-ratio, and the posterior probability of a positive and negative coefficient value were obtained. These quantities enable interpretation of covariate effects along with their uncertainty and not just by using point estimates.

The sensitivity analysis examined training information, administrative censoring, predictor definitions, missingness, baseline-hazard intervals and prior scales. Explanation analysis compared posterior contribution rankings with exact linear SHAP rankings, and checked their stability across sampling budgets and seeds. These analyses retained the same separation between fitting, diagnostic acceptance and interpretation.

\begin{figure}[p]
  \centering
  \begingroup
  \singlespacing
  \begin{tikzpicture}[
    x=1cm,y=1cm,
    box/.style={draw=black!65,rounded corners=2pt,line width=0.6pt,
      align=center,font=\fontsize{10}{12}\selectfont,inner sep=5pt,execute at begin node={\hyphenpenalty=10000},text width=8.2cm},
    attempt/.style={box,fill=blue!7},
    gate/.style={box,fill=violet!8},
    arrow/.style={-{Stealth[length=2mm]},draw=black!70,line width=0.65pt},
    edge label/.style={font=\small,fill=white,inner sep=1pt}
  ]
    \node[attempt,minimum height=1.2cm] (initial) at (10,0)
      {\textbf{Initial Bayesian fit}\\
       Four chains; 2,000 warm-up and 2,000 retained draws per chain};
    \node[gate,minimum height=2.2cm] (gateone) at (10,-2.3)
      {\textbf{All diagnostic criteria satisfied?}\\
       Zero divergences and tree-depth hits\\
       $\min\mathrm{BFMI}\geq0.30$; $\max\widehat R<1.01$; $\min\mathrm{ESS}_{\rm bulk}>400$\\
       Finite aggregate BFMI, $\widehat R$ and bulk ESS};
    \node[attempt,minimum height=1.3cm] (remedyone) at (10,-4.6)
      {\textbf{First specified refit}\\
       Target acceptance 0.99; 4,000 warm-up iterations; changed seed};
    \node[gate,minimum height=0.9cm] (gatetwo) at (10,-6.25)
      {\textbf{Same diagnostic gate satisfied?}};
    \node[attempt,minimum height=1.3cm] (remedytwo) at (10,-7.9)
      {\textbf{Second specified refit}\\
       Retain first-refit controls; non-centred baseline log-hazards\\
       Identical statistical prior and likelihood};
    \node[gate,minimum height=0.9cm] (gatethree) at (10,-9.65)
      {\textbf{Same diagnostic gate satisfied?}};
    \node[box,fill=orange!12,minimum height=1.3cm] (failed) at (10,-11.4)
      {\textbf{Record diagnostic status}\\
       Record diagnostic summaries;\\
       paired inference uses accepted fits};
    \node[box,fill=teal!8,text width=3.25cm,minimum height=2.5cm] (accepted) at (2,-6.25)
      {\textbf{Retain the first passing attempt}\\[0.3em]
       Check prediction integrity, then use eligible predictions and posterior summaries};

    \draw[arrow] (initial.south) -- (gateone.north);
    \draw[arrow] (gateone.south) -- node[edge label,right]{Fail} (remedyone.north);
    \draw[arrow] (remedyone.south) -- (gatetwo.north);
    \draw[arrow] (gatetwo.south) -- node[edge label,right]{Fail} (remedytwo.north);
    \draw[arrow] (remedytwo.south) -- (gatethree.north);
    \draw[arrow] (gatethree.south) -- node[edge label,right]{Fail} (failed.north);
    \draw[draw=black!70,line width=0.65pt] (4.65,-2.3) -- (4.65,-9.65);
    \foreach \test in {gateone,gatetwo,gatethree}
      \draw[draw=black!70,line width=0.65pt] (\test.west) --
        node[edge label,above]{Pass} (4.65,0 |- \test.west);
    \draw[arrow] (4.65,-6.25) -- (accepted.east);
  \end{tikzpicture}
  \endgroup
  \caption[Bayesian sampling and diagnostic assessment.]{Bayesian sampling and diagnostic assessment. Each fit is assessed against the same criteria after the initial attempt and each specified refit. The first accepted attempt supplies posterior and predictive summaries. The horseshoe model starts at target acceptance 0.99. Summary precision and the roles of tail ESS and MCSE are specified in the methods.}
  \label{fig:diagnostic-decisions}
\end{figure}


\section{Data set and Cohort Preparation.}\label{data-set-and-cohort-preparation.}

\subsection{Healthcare Cohort Selection}\label{healthcare-cohort-selection}

Three publicly available healthcare survival data sets were used in this empirical investigation. They are: GBSG2, METABRIC and WHAS500. They all have varying clinical information, though they are all time-to-event modelling. The importance to have several independent data sets is based on the fact that the survival models can be severely affected by the sample size, the structure of predictors, censoring and clinical situations. The recent literature on survival modelling emphasizes the significance of assessing survival prediction models using a wide range of data sets \cite{ref33,ref44}.

GBSG2 (German Breast Cancer Study Group 2) is a study cohort of breast-cancer survival. In the ultimate implementation, the sample size is 686 and the predictors are 8. The predictor set includes the age, tumour size (tsize), positive nodes (pnodes), progesterone receptor measurement (progrec), oestrogen receptor measurement (estrec), hormonal therapy (horTh), menopausal status (menostat) and tumour grade (tgrade). Recurrence-free survival is considered the endpoint and the time of survival is in days. The implementation makes 299 observed and 387 censored \cite{schumacher1994,polsterl2020}.

A second data set on breast cancer is METABRIC (Molecular Taxonomy of Breast Cancer International Consortium). The reconstructed confirmatory data set of 1980 observations and 8 predictors that include age, Nottingham Prognostic Index (NPI), cellularity, ER status, PR status, HER2 status, menopausal status and PAM50/claudin-low subtype is the final reconstructed data set. The total survival (in months) is the outcome variable. The data set consists of 1144 observed events, and 836 censored observations. METABRIC has been heavily used in recent breast cancer survival modelling investigations including using clinical and molecular information for prognostic modelling \cite{curtis2012,pereira2016}.

WHAS500 (Worcester Heart Attack Study) does not offer cancer-related survival data, but cardiovascular. The implementation includes 500 observations and 9 major predictors such as age, gender, heart rate, systolic blood pressure, diastolic blood pressure, body mass index, cardiovascular-disease history, order of myocardial infarction and type of myocardial infarction. All cause mortality, in days is the outcome variable. There are 215 observed events and 285 censored events in the data set. WHAS500 is a common criterion to assess survival learning techniques \cite{hosmer2008,polsterl2020}.

\subsection{Data Reconstruction and acquisition.}\label{data-reconstruction-and-acquisition.}

These data were considered as secondary, publicly available, survival data. The confirmatory implementation obtained the GBSG2 and WHAS500 cohorts with the default scikit-survival dataset loaders. METABRIC was rebuilt using an established cohort-specific process to project the source variables into a frozen primary analysis table. This is significant since the completed modelling dataset is determined by the confirmatory pipeline that has been validated as opposed to the initial exploratory CSV representation.

The reconstruction procedure made a clear distinction between predictors and survival outcomes: GBSG2 time and cens, WHAS500 lenfol and fstat, and METABRIC time and event. Source fields were mapped to explicit binary event indicators according to their documented coding. In particular, the GBSG2 field named cens supplies the observed-event indicator. The cohort adapters standardised the representation to time and event, enabling a common survival-analysis interface across the three cohorts.

\subsection{Predictor Preparation}\label{predictor-preparation}

The implementation identifies numerical, ordinal, and nominal categorical predictors instead of making a single transformation of all variables. In the case of GBSG2, the tumour grade was considered as an ordered variable with grades I, II and III being deterministically represented by 0, 1 and 2. Numerical predictors were median-imputed and standardised using training statistics. Ordered values were encoded before standardisation, with missing or unknown values represented as $-1$.

Nominal categorical variables were considered independently. Missing categorical values were coded as an explicit missing category, and then coded as one-hot with drop=first and handle\_unknown=ignore. This is to make sure that categorical levels missing in a given validation fold will not lead to transformation failure.

After transformation with the output of the fitted preprocessor (feature-names), names of features were stored. This is relevant to future interpretation and explainability of Bayesian coefficients since it is necessary that the transformed positions of predictors can be traced back to their original variables.

\subsection{Cohort-Specific Preprocessing and Control of Leakage.}\label{cohort-specific-preprocessing-and-control-of-leakage.}

The analysis does not fit one global preprocessing transformation to the complete dataset; instead, preprocessing is fitted on each outer-training partition and applied unchanged to its validation partition. The pipeline verifies training and validation dimensional consistency, finite transformed values and the lack of train validation overlap.

This ensures that validation observations do not determine imputation statistics, scaling or categorical encoding in the model. This is especially relevant in the survival prediction because it may add optimistic estimates of the performance in case pre-processing and validation choices are made \cite{ref32,ref34}.

\begin{longtable}[]{@{}
  >{\raggedright\arraybackslash}p{(\columnwidth - 12\tabcolsep) * \real{0.1600}}
  >{\raggedright\arraybackslash}p{(\columnwidth - 12\tabcolsep) * \real{0.2093}}
  >{\raggedright\arraybackslash}p{(\columnwidth - 12\tabcolsep) * \real{0.1195}}
  >{\raggedright\arraybackslash}p{(\columnwidth - 12\tabcolsep) * \real{0.1180}}
  >{\raggedright\arraybackslash}p{(\columnwidth - 12\tabcolsep) * \real{0.1416}}
  >{\raggedright\arraybackslash}p{(\columnwidth - 12\tabcolsep) * \real{0.1227}}
  >{\raggedright\arraybackslash}p{(\columnwidth - 12\tabcolsep) * \real{0.1290}}@{}}
\caption{Canonical study cohorts, predictor counts, follow-up scales, and event/censoring totals}\label{_Toc245465016}\\
\toprule\noalign{}
\begin{minipage}[b]{\linewidth}\raggedright
Cohort
\end{minipage} & \begin{minipage}[b]{\linewidth}\raggedright
Clinical domain
\end{minipage} & \begin{minipage}[b]{\linewidth}\raggedright
Records (n)
\end{minipage} & \begin{minipage}[b]{\linewidth}\raggedright
Source variables
\end{minipage} & \begin{minipage}[b]{\linewidth}\raggedright
Time scale
\end{minipage} & \begin{minipage}[b]{\linewidth}\raggedright
Events (n)
\end{minipage} & \begin{minipage}[b]{\linewidth}\raggedright
Censored (n)
\end{minipage} \\
\midrule\noalign{}
\endhead
\bottomrule\noalign{}
\endlastfoot
GBSG2 & Breast cancer & 686 & 8 & Days & 299 & 387 \\
METABRIC & Breast cancer & 1,980 & 8 & Months & 1,144 & 836 \\
WHAS500 & Cardiovascular / myocardial infarction & 500 & 9 & Days & 215 & 285 \\
\end{longtable}

Table~\ref{_Toc245465016} reports the cohort definitions used in the analysis. Predictor counts refer to source variables before categorical encoding; the primary METABRIC cohort contains 1,980 observations and eight source predictors.

\section{Survival Outcome Definition}\label{survival-outcome-definition}

\subsection{Time-to-Event Representation}\label{time-to-event-representation}

In this study, the outcome is a time-to-event variable. For individual $i$, let the event time be:

\[T_{i}\]

Whilst the three data sets have different clinical outcomes and time units, they are all based on survival modelling. GBSG2 models are recurrence free survival, WHAS500 models are all cause mortality after a cardiovascular event and METABRIC models are overall survival.

The observed follow-up records either the event time or the censoring time. Recent studies on survival prediction emphasize the importance of preserving time varying and time invariant information in time-to-event modelling to enable individual level predictions to be made \cite{ref14,ref45}.

\subsection{Event Indicator}\label{event-indicator}

The observed event indicator of each individual is determined to be:
\begin{equation}
    \delta_{i} = \left\{ \begin{matrix}
    1, & \text{if~the~event~is~observed}, \\
    0, & \text{if~the~observation~is~right~censored}.
    \end{matrix} \right.\ \
\end{equation}


The cohort-specific outcome fields are converted to a unified binary event indicator. GBSG2 uses cens as the observed-event indicator, WHAS500 uses fstat, and METABRIC provides an event field.

\subsection{Right Censoring}\label{right-censoring}

Let \emph{C\textsubscript{i}} denote the potential censoring time. The recorded follow-up is the earlier of the event time and censoring time:

\begin{equation}
    T_i^{\mathrm{obs}}=\min(T_i,C_i).
\end{equation}



If the event happens before the censoring time occurs, we have an event (\emph{T\textsubscript{i}} \ensuremath{\leq} \emph{C\textsubscript{i}} ). The event variable will be 1, and the observed time will be \emph{T\textsubscript{i}}. If the event does not occur before the censoring time, we have no event (event variable = 0), and the observed time will be \emph{C\textsubscript{i}}. Having no event does not mean the event will not occur, just that it will not be observed within the follow-up period. Treating censored observations as negative binary outcomes would discard their known event-free follow-up and can bias the analysis \cite{ref14,ref45}.

\subsection{Survival and Hazard Functions}\label{survival-and-hazard-functions}

The survival function can be defined as
\begin{equation}
    S(t) = P(T > t),
\end{equation}


The survival function gives the probability of remaining event-free beyond time t. The corresponding hazard function is defined as:

\begin{equation}
    h(t) = \lim_{\Delta t \rightarrow 0}\frac{P(t \leq T < t + \Delta t \mid T \geq t)}{\Delta t}.
\end{equation}



The survival function provides us with the probability of surviving at a given time. The hazard provides us with the rate of the event occurring at a specific time, given that the subject is alive. These can be used to help with the prediction problem in the thesis as they take into account the time of the event, rather than binary classification. Modern survival prediction models explicitly take this into account \cite{ref14,ref26}.

Evaluation metrics such as time-dependent AUC and Brier prediction error take into account the time aspect and censoring \cite{ref26,ref47}.

\section{Train/Test Separation and Validation Design}\label{traintest-separation-and-validation-design}

The validation strategy was designed to avoid overfitting and to ensure the 3 healthcare cohorts had the time-to-event structure. Censoring, discrimination, prediction error and calibration need to be considered when predicting survival. Therefore, the evaluation strategy didn't rely on a single measure of performance. A review of methods of evaluating survival prediction highlights that discrimination, calibration and prediction error complement each other \cite{ref26,ref27}.

The adopted pipeline created an outer-validation structure that included 5 repeats and 5 folds per cohort with resulting 25 outer-fold definitions per cohort and 75 outer-fold definitions across GBSG2, METABRIC and WHAS500. Partitions were stratified by event status and, where group sizes permitted, ranked follow-up-time bins. Repeat-specific seeds were derived from the base seed. The explicit training and validation indices were present in each outer fold. This resulting structure gave a prespecified division between observations to be used in developing the model and those to be used in held-out evaluation.

An important methodological aspect was fold-local preprocessing. In the case of every outer fold, the preprocessing state was obtained based on the respective training partition and then applied to the validation partition. This makes sure that information in the held-out validation observations does not affect feature transformation parameters. Training and validation indices, transformed predictor matrices, survival times, event indicators, preprocessing state and supported reporting horizons were stored as fold-specific objects in the fold-preparation routine. As a result, the validation observations were not part of the training transformation process but served as an assessment set.

The modelling pipeline applied the same outer partitions to Cox proportional hazards regression, Random Survival Forest, DeepSurv and the three Bayesian variants. RSF and DeepSurv used three-fold inner cross-validation within each outer-training partition to select the configuration with the lowest mean inner IBS. Preprocessing was fitted separately within each inner-training partition. The selected configuration was then refitted within the outer-training data. Each patient appeared in outer validation once per repeat. Metrics were calculated within folds and summarised across repeats, rather than pooling repeated predictions as independent patient observations.

\section{Baseline Survival Models}\label{baseline-survival-models}

Three different survival model families were selected to compare with the Bayesian survival models. Cox proportional hazards (CoxPH), Random Survival Forest (RSF) and DeepSurv were implemented. This was done to provide different reference models. The performance of these models can be compared to each other based on the amount of events, censoring, structure of predictors, linearity, interactions and evaluation metrics. This can be done using comparative survival analysis \cite{ref33,ref37,ref44}.

\subsection{Cox Proportional Hazards}\label{cox-proportional-hazards}

As the traditional statistical baseline, the Cox proportional hazards model was applied. Its hazard formulation is

\begin{equation}
    h_{i}(t) = h_{0}(t)\exp\left( \mathbf{x}_{i}^{T}\mathbf{\beta} \right),
\end{equation}


Here, $h_0(t)$ denotes the baseline hazard, $\mathbf{x}_i$ the predictor vector for subject $i$, and $\boldsymbol{\beta}$ the coefficient vector.

The final CoxPH model was fitted using only the outer training fold with a fixed penalizer of 1x10\textsuperscript{-}\textsuperscript{4}. No hyperparameter tuning was performed for the outcome driven CoxPH parameters. Numerical variables and categorical variables with encoding from the outer fold were passed into lifelines. The fitted baseline survival and linear predictor provide absolute survival probabilities for the holdout records. The coefficients are exponentiated to provide the relative hazard ratios. CoxPH provides an appropriate reference model as the Bayesian model has the same proportional hazards structure. This assumes the hazard ratios stay constant over time for the covariates and the linear predictor. If there are time varying effects or non-linear relationships, CoxPH may not be suitable. RSF provides a non-linear reference model. DeepSurv provides a less linear model but still has proportional hazards.

Rank-transformed Schoenfeld-residual tests assessed proportional-hazards departures, with Holm adjustment within each training-fit covariate family \cite{schoenfeld1982,holm1979}.

Recent simulations have shown that neither model is always better. CoxPH can provide as strong performance in the presence of linear effects and little information on the event. RSF can benefit from situations with non-linearities and interactions \cite{ref37}.

\subsection{Random Survival Forest}\label{random-survival-forest}

The RSF model had 300 trees, a minimum split size of 10, no maximum depth and one computation thread. Within each outer-training fold, three-fold inner cross-validation compared leaf sizes of 5 or 15 and candidate-feature sampling using either the square root of the predictor count or all predictors. The selected model was refitted using the outer-training data. The risk score is the ensemble prediction, and survival probabilities are calculated using the estimated survival functions from the forest. The tree ensemble can represent nonlinear relationships and interactions without using the linear predictor from Cox.

The random survival forest provides a nonlinear comparator for CoxPH and the Bayesian proportional-hazards models. Its results help assess whether the additional Bayesian structure is associated with improved predictive performance in these cohorts. RSF has been used in high dimensional censored survival data \cite{ref51}.

RSF provides a predictive comparator, while the explanation analysis focuses on the Bayesian linear log-risk scores and the CoxPH reference.

\subsection{DeepSurv}\label{deepsurv}

The DeepSurv model used three-fold inner cross-validation to select among four configurations. These combined hidden widths $[32]$ or $[32,16]$ with either no dropout and $L_2=10^{-4}$, or dropout 0.2 and $L_2=10^{-3}$. All configurations used ReLU activation, batch normalisation, the Adam optimiser with learning rate 0.001, batch size 64 and a maximum of 100 epochs. An event-stratified 20\% subset of the current training data controlled early stopping, with patience 15 and minimum improvement $10^{-4}$. The log-risk score was combined with the Breslow baseline from the training fold to predict holdout survival probabilities. DeepSurv allows for nonlinear predictor effects but still uses the Cox proportional hazards objective function \cite{ref4}.

Katzman et al. have shown that DeepSurv performs well on simulated and clinical survival data. It can be seen as a neural version of the Cox survival model for personal risk prediction \cite{ref4}. The model provides a very different computational reference but still has a familiar survival risk prediction structure. The three baselines are used jointly to denote conventional statistical modelling (CoxPH), the nonlinear ensemble learning (RSF), and neural survival modelling (DeepSurv). This combination offers a significant methodological reference to assess the Bayesian piecewise-exponential proportional-hazards models.

The three baseline families were evaluated on all 25 outer folds in each cohort, using the same held-out records as the Bayesian models.

\section{Bayesian Piecewise-Exponential Proportional-Hazards Model}\label{bayesian-piecewise-exponential-proportional-hazards-model}

The main Bayesian aspect of the study was developed as a completely Bayesian piecewise-exponential proportional-hazards (PH) model. The implemented Bayesian model, as opposed to the Cox proportional-hazards baseline comparator, which uses a partial likelihood, has an explicit representation of the baseline hazard. The model specification was frozen prior to confirmatory inference and applicable a piecewise-constant baseline hazard broken down into six intervals determined from training event-time quantiles. An extended Poisson representation of the survival likelihood was also used in the implementation as the time-to-event observations could be encoded as interval-specific exposure and event records.

The proportional-hazards formulation of an individual, i, is:
\begin{equation}
    h_{i}(t) = h_{0}(t)\exp\left( \mathbf{x}_{i}^{T}\mathbf{\beta} \right),
\end{equation}


Here, $h_i(t)$ is the hazard at time $t$, $\mathbf{x}_i$ is the predictor vector, and $\boldsymbol{\beta}$ is the coefficient vector. The baseline hazard is denoted as

\begin{equation}
    h_{0}(t) = \lambda_{j},\quad\quad t \in I_{j},\quad j = 1,\ldots,6,
\end{equation}


Here, $I_j$ is the $j$th interval and $\lambda_j$ its hazard rate. Therefore, the model enables the baseline hazard to vary across intervals, but with a proportional covariate effect structure across intervals.

The enlarged representation converts the follow-up of every subject into interval data. In interval number, an event-count variable is the event that happened during that interval and the corresponding exposure/offset is the time that the subject contributed. The resulting Poisson formulation can be conceptualized as
\begin{equation}
    y_{ij} \sim \text{Poisson}\left( E_{ij}\lambda_{j}\exp\left( \mathbf{x}_{i}^{T}\mathbf{\beta} \right) \right),
\end{equation}

Here, $y_{ij}$ is the interval event count and $E_{ij}$ the exposure. This provides a computationally friendly Bayesian inference approach. Piecewise exponential Bayesian models have also been used to model time-to-event processes with predefined time intervals. A more general Bayesian survival methodology has been used to include uncertainty in the posterior \cite{laird1981,friedman1982,ref22}.

Within each training fold, time was divided by the median observed event time. The six intervals used equally spaced quantiles of scaled training event times, with zero as the first boundary and maximum scaled training follow-up plus one as the final boundary. Equal-width intervals were used if tied quantiles produced too few distinct intervals. The shared baseline prior was $\log\lambda_j\sim N(-2,2^2)$ on this scaled time axis. The likelihood equations use this scaled time; converting the hazard back to the original time unit divides its rate by the training median event time. Prediction beyond the final boundary continued the last rate. Survival probabilities were obtained for each posterior draw and then averaged; the ranking score used the posterior-mean linear predictor.

The posterior distribution is thus:

    \begin{equation}
    p\left( \mathbf{\beta},\mathbf{\lambda} \mid D \right) \propto p\left( D \mid \mathbf{\beta},\mathbf{\lambda} \right)p\left( \mathbf{\beta} \right)p\left( \mathbf{\lambda} \right),
\end{equation}


where \(D\) is the observed survival data. The use of posterior samples gives distributions and not just point estimates, allowing the determination of posterior means, 95\% credible intervals, hazard-ratio distributions, and directional probabilities of individual predictors.

There were three alternative coefficient-prior structures used to apply the model on GBSG2, METABRIC and WHAS500. These Bayesian variants also maintained the six-interval baseline specification and full-likelihood formulation, so that differences between Bayesian models could be largely due to the prior structure rather than to alterations to the survival likelihood. The application thus offers a regulated framework of researching into the impact of prior presumptions on prediction as well as posterior clarification.

\section{Bayesian Prior Specification}\label{bayesian-prior-specification}

Prior specification was also an important aspect of this thesis since the objective of this research was to determine whether various assumptions around coefficient shrinkage have any effect on inference, predictions and explainability. The Normal, Horseshoe, and Spike-and-Slab priors were compared using the same piecewise-exponential likelihood and six baseline intervals, allowing differences in posterior behaviour to be examined under a common model structure. Various shrinkage and selection priors were employed and this is consistent with the recent Bayesian survival analyses that have employed shrinkage and selection structures to control complexities of the coefficients and to determine predictors useful to survival \cite{ref42,ref43}.

\subsection{Normal Prior}\label{normal-prior}

The reference Bayesian specification assigned
\begin{equation}
    \beta_{k} \sim N(0,1)
\end{equation}


to each regression coefficient. The zero-centred prior offers a moderate level of regularisation with both positive and negative effects being present. It thus is the relatively traditional Bayesian reference with which the more structured shrinkage methods may be compared. Normal-prior model does not explicitly select variables: instead, the uncertainty in the posterior is used to find out how much individual coefficients can be distinguished against zero.

\subsection{Horseshoe Prior}\label{horseshoe-prior}

The Horseshoe specification uses a regularised global-local shrinkage structure:
\begin{align}
\tau&\sim t^+_2(0,0.1), & q_k&\sim t^+_2(0,1),\\
c^2&\sim\operatorname{InvGamma}(2,8), &
\widetilde q_k^2&=\frac{c^2q_k^2}{c^2+\tau^2q_k^2},\\
z_k&\sim N(0,1), & \beta_k&=z_k\tau\widetilde q_k.
\end{align}
The global scale $\tau$ controls overall shrinkage, local scales $q_k$ allow coefficient-specific shrinkage, and the slab parameter $c^2$ limits extreme coefficient scales. Here, $t^+_2$ denotes a half-Student-$t$ distribution with two degrees of freedom. These half-$t_2$ scales specify the variant used in this study \cite{carvalho2010,piironen2017}. Empirical suitability was assessed using the same diagnostic acceptance criteria as for the other priors.

\subsection{Spike-and-Slab Prior}\label{spike-and-slab-prior}

The continuous Spike-and-Slab used a variance mixture:
\begin{align}
\gamma_k&\sim\operatorname{Beta}(1,1),\\
\sigma_k^2&=(1-\gamma_k)(0.05)^2+\gamma_k(2.0)^2,\\
z_k&\sim N(0,1),\qquad \beta_k=z_k\sigma_k.
\end{align}
The two variances were mixed using a continuous weight. The model did not have a discrete feature-membership indicator. The posterior mean of $\gamma_k$ therefore describes variance allocation rather than a probability of variable inclusion.

\subsection{Rationale for Comparing Prior Families.}\label{rationale-for-comparing-prior-families.}

The three priors were then considered as competing assumptions under the same likelihood framework as opposed to an argument that one prior is universally superior. Their posterior behaviour, predictive performance, and feature-level uncertainty is assessed in the following analysis. This is especially crucial since the smaller the posterior coefficient does not imply a better model; shrinkage needs to be viewed in conjunction with predictive discrimination, prediction error, calibration and credible uncertainty.

\section{MCMC Inference and Computational Configuration}\label{mcmc-inference-and-computational-configuration}

Bayesian posterior inference was done using Hamiltonian Monte Carlo (HMC) with the No-U-Turn Sampler (NUTS). This can be used for Bayesian models with multiple continuous parameters as it uses the gradient information of the posterior to better navigate through parameter space than most traditional MCMC random-walk samplers. More modern approaches have demonstrated that HMC can be quite helpful in computing the Bayesian way and it is necessary to learn how the sampler works \cite{ref55}.

The confirmatory MCMC configuration was prespecified and applied to the Bayesian fits, with prior-specific target acceptance probabilities as reported below. Four independent chains were sampled in each Bayesian model, and 2,000 warm-up iterations and 2,000 retained draws were used in each chain. Therefore, every model that was accepted produced 8,000 retained posterior draws on the four chains. Deterministic seeds for cohort, repeat, fold, model and attempt were derived from base seed 260817. The target acceptance probability of the Normal and Spike-and-Slab models was 0.95 and 0.99 of the Horseshoe model, respectively, showing the extra computational cost of hierarchical global-local shrinkage. Bayesian inference used the six-interval piecewise-exponential baseline representation and the training data which were specific to the cohort.

This analysis was divided into three sections, posterior sampling, diagnostics and posterior predictive analysis. This is significant because one cannot assume that by simply having an MCMC algorithm run correctly the draws will be adequate to sample the target posterior completely. Recent Bayesian guidelines workflows concentrate on the computational validation and posterior interpretation \cite{ref53}.

Controlled diagnostic-recovery procedures of problematic posterior fits were also implemented. These operations changed as necessary sampling controls or parameterisation, but held the original likelihood, predictors six baseline intervals and diagnostic acceptance criteria. Recovery thus was an effort to enhance computational exploration and not a process of loosening model-quality criteria.

Prior predictive and posterior predictive checks each used 200 draws, providing complementary assessments of the prior implications and model fit. Bayesian survival analysis emphasizes the importance of computational inference and posterior diagnostics \cite{ref22,ref53}.

\begin{longtable}[]{@{}
  >{\raggedright\arraybackslash}p{(\columnwidth - 2\tabcolsep) * \real{0.5516}}
  >{\raggedright\arraybackslash}p{(\columnwidth - 2\tabcolsep) * \real{0.4484}}@{}}
\caption{Bayesian sampler, chain, warm-up, draw, seed, and predictive-check configuration}\label{_Toc245465017}\\
\toprule\noalign{}
\begin{minipage}[b]{\linewidth}\raggedright
Configuration
\end{minipage} & \begin{minipage}[b]{\linewidth}\raggedright
Implemented specification
\end{minipage} \\
\midrule\noalign{}
\endhead
\bottomrule\noalign{}
\endlastfoot
Bayesian sampler & PyMC NUTS \\
Chains & 4 \\
Warm-up/tuning & 2,000 \\
Retained draws & 2,000 per chain \\
Total retained draws & 8,000/model \\
Base random seed & 260817 \\
Normal target acceptance & 0.95 \\
Horseshoe target acceptance & 0.99 \\
Spike-and-Slab target acceptance & 0.95 \\
Prior predictive draws & 200 \\
Posterior predictive draws & 200 \\
Baseline intervals & 6 \\
\end{longtable}



\section{MCMC Diagnostic Framework and Model Acceptance}\label{mcmc-diagnostic-framework-and-model-acceptance}

MCMC diagnostics determined whether each posterior fit was accepted for interpretation. This is important as the posterior summaries can be misleading if the chains have not fully sampled the target distribution. Recent literature in Bayesian workflow indicates that computational diagnostics, model checking and sensitivity are a consideration of the Bayesian analysis \cite{ref53}.

The following diagnostics were performed to check the posterior: divergences, tree depth saturation, Bayesian fraction of missing information (BFMI), bulk and tail effective sample size (ESS), rank normalised R-hat, Monte Carlo standard error (MCSE) and posterior finiteness.

First, the divergences had to be zero. Problematic posterior geometry was represented by divergent NUTS transitions and thus was a failure of the strict acceptance gate. This was particularly important for the hierarchical Horseshoe specification, whose additional global-local parameters can generate challenging posterior geometries.

Second, the maximum tree-depth hits needed to be zero. The saturation of tree-depths may hint that the sampler has constantly attained a limit of its computational trajectory without fully exploring it.

Third, BFMI was measured in chains that had minimum required value of 0.30. BFMI was an energy-based diagnostic of whether the sampler was exploring adequately the posterior momentum distribution.

The bulk and tail ESS were checked. The minimum bulk ESS had to be more than 400; tail ESS was recorded as a supplementary diagnostic. The sample size is taken into consideration with the autocorrelation of the MCMC chains. This describes the information in autocorrelated draws relative to independent posterior sampling. The recent MCMC post-processing literature points out that the amount of draws is not to be confused with the amount of information being utilized \cite{ref39,ref54}.

The highest possible rank normalised R-hat value of all parameters must be less than 1.01. This gives a more accurate measure of agreement between the chains in case the posterior distributions are skewed or heavy tailed \cite{ref39}.

Sixth, MCSE was computed and stored as another indicator of Monte Carlo uncertainty. Though MCSE was not directly employed as the main binary exclusion criterion, it offered an extra quantitative measure of the accuracy of posterior Monte Carlo estimates.

Posterior predictive checks (PPC) are conducted to verify whether the data that is replicated by the model fitted resemble the observed data. This is independent of the MCMC convergence.

The strict rule of acceptance so resulting was

\begin{equation}
    \begin{aligned}
    \mathrm{Accept} = \mathbb{I}\{ & D = 0,\ T = 0,\ BFMI_{\min} \geq 0.30, \\
     & ESS_{bulk,min} > 400, \\
     & {\widehat{R}}_{\max} < 1.01,\ F = 1\}.
    \end{aligned}
\end{equation}



Here, $D$ is the divergence count, $T$ is the tree-depth-hit count and $F$ indicates finite BFMI, R-hat and bulk-ESS aggregates.

The gate used the recorded summary precision: R-hat rounded to two decimal places and ESS to integers. BFMI, R-hat and bulk-ESS aggregates had to be finite. Following an initial failure, the first specified refit used 4,000 warm-up iterations and target acceptance 0.99. The second used the same controls and a non-centred representation of the identical baseline prior. The first passing attempt supplied the posterior and predictive summaries. This procedure was applied to every Bayesian model--fold fit.

\begin{longtable}[]{@{}
  >{\raggedright\arraybackslash}p{(\columnwidth - 2\tabcolsep) * \real{0.4811}}
  >{\raggedright\arraybackslash}p{(\columnwidth - 2\tabcolsep) * \real{0.5189}}@{}}
\caption{Prespecified acceptance criteria for Bayesian sampling diagnostics}\label{posterior-predictive-checking}\\
\toprule\noalign{}
\begin{minipage}[b]{\linewidth}\raggedright
Diagnostic
\end{minipage} & \begin{minipage}[b]{\linewidth}\raggedright
Frozen acceptance criterion
\end{minipage} \\
\midrule\noalign{}
\endhead
\bottomrule\noalign{}
\endlastfoot
Divergences & \textbf{0} \\
Maximum tree-depth hits & \textbf{0} \\
Minimum BFMI & \textbf{\ensuremath{\geq} 0.30} \\
Minimum bulk ESS & \textbf{\textgreater{} 400} \\
Tail ESS and MCSE & Recorded as supplementary summaries \\
Maximum R-hat & \textbf{\textless{} 1.01} \\
Finite diagnostic aggregates & \textbf{Yes} \\
Diagnostic bypass & \textbf{No} \\
Fit acceptance & Assessed separately for each outer fold \\
\end{longtable}

\section{Posterior Predictive Checking}\label{posterior-predictive-checking-1}

Posterior predictive checks (PPC) are conducted to verify whether the data that is replicated by the model fitted resemble the observed data. This is independent of the MCMC convergence.

For each accepted Bayesian fit, posterior predictive checks used up to 200 evenly spaced retained draws. Poisson event counts were replicated for every positive-exposure likelihood row and summed both overall and by baseline interval. Observed counts were compared with the replicated means, medians and central 95\% intervals. A fitted-survival check also compared training-average predicted survival with Kaplan--Meier estimates at training event-time quartiles.

This maintains the methodological order:
\[\text{MCMC sampling}\rightarrow\text{diagnostic acceptance}\rightarrow\text{posterior predictive checking}.\]
Posterior predictive checks describe the fitted model and complement the separate held-out assessment of predictive performance.

\section{Predictive Evaluation Framework}\label{predictive-evaluation-framework}

The predictive evaluation model was created to see how well the model discriminates and predicts accurately the censored time-to-event results. Survival prediction is the prediction of the likelihood of a certain event happening and at what time. Thus, there is no single measure that can be applied to adequately describe the model performance. Contemporary methodological literature recommends discrimination, prediction error and calibration to be taken into account. C-index informs us on the extent to which the model is able to rank patients based on the relative risk of events. Brier score informs us of the extent to which the forecasted survival is not similar to the actual. Time-dependent AUC informs us on how the model can differentiate at a certain follow-up time. The agreement between the predicted and observed probability of an event is measured by calibration \cite{ref26,ref27,ref59}.

The likelihood assumes censoring is independent of the event time conditional on the included predictors. The marginal IPCW estimators additionally assume independent censoring over their supported evaluation window.

The primary measure of probability prediction error was IPCW Integrated Brier Score (IBS):
\begin{equation}
\mathrm{IBS}=\frac{1}{t_U-t_L}\int_{t_L}^{t_U}\widehat{BS}(t)\,dt.
\end{equation}
The Brier score is the inverse-probability-of-censoring-weighted squared distance between the observed survival indicator and predicted survival probability at time $t$. The reporting horizons were the 25th, 50th and 75th percentiles of observed training event times; IBS used 25 equally spaced grid points between the first and last horizons. The implementation used scikit-survival. Lower IBS indicates smaller prediction error.

Uno's concordance was the primary discrimination measure. Harrell's concordance and cumulative/dynamic AUC provided additional summaries, with larger values indicating better discrimination. The training Kaplan--Meier estimate of censoring survival had to exceed 0.10 over the evaluation grid. For IPCW IBS, Uno concordance and AUC, validation follow-up beyond the final training horizon was administratively censored immediately after that horizon, preserving event-free status throughout the supported window. Harrell concordance and calibration used observed validation follow-up.

Calibration assessed agreement between predicted and observed event probabilities. At each horizon, with $z_i(t)=\log[-\log\widehat S_i(t)]$, an IPCW complementary-log--log model estimated
\begin{equation}
\operatorname{cloglog}\{\Pr(T_i\leq t)\}=\alpha_t+\beta_t z_i(t).
\end{equation}
The ideal joint intercept and slope are zero and one. Calibration-in-the-large was estimated separately as the intercept with slope fixed at one. An observed event received weight $1/\widehat G(T_i^-)$, known event-free status beyond the horizon received $1/\widehat G(t)$, and censoring by the horizon received weight zero. Calibration intervals used 500 event-stratified bootstrap replicates, resampling training observations to re-estimate censoring and validation observations with their predictions. At least 80\% of replicates, and at least 20, had to be estimable. Five equal-count prediction groups provided descriptive calibration curves.

Model contrasts used paired outcomes from identical cohort, repeat and fold partitions. For 25 paired metric differences with sample variance $s_d^2$, the corrected repeated-cross-validation standard error was
\begin{equation}
\operatorname{SE}_{\mathrm{corr}}=\sqrt{\left(\frac{1}{25}+\frac{1}{4}\right)s_d^2}.
\end{equation}
Nominal 95\% $t$ intervals used 24 degrees of freedom \cite{nadeau2003}. A contrast required all 25 paired accepted outcomes. Negative IBS differences favour the first named model. The nine-test primary family comprised three pairwise prior contrasts in each cohort, with Holm adjustment specified at family-wise level 0.05 \cite{holm1979}. Comparisons with CoxPH, RSF and DeepSurv were secondary. A paired event-stratified bootstrap with 1,000 replicates assessed sensitivity to the corrected interval method. These estimates concern predictive differences; no equivalence margin was specified.

\section{Bayesian Explainability Framework}\label{bayesian-explainability-framework}

The explainability framework was built on the probabilistic form of the Bayesian survival models as opposed to solely depending on the post-hoc black-box explanation techniques. Bayesian inference is a direct characterization of the uncertainty in the parameters in the form of posterior distributions, making it possible to interpret the model by including the estimated magnitude of effects and the uncertainty. This is especially suitable when one wants to conduct a survival analysis since the posterior hazard ratios can be interpreted in a probabilistic sense and not just by point estimates.

The posterior coefficient of predictor k of each retained Bayesian model was modeled as.
\begin{equation}
    \beta_{k} \mid D,
\end{equation}


where \(D\) is the observed survival data. The hazard ratio was found to be:
\begin{equation}
    HR_{k} = \exp\left( \beta_{k} \right)
\end{equation}


The posterior summaries cover accepted Bayesian fits across the repeated outer validation. First-fold examples provide a readable illustration of the feature-level estimates; the repeated-fold explanation analysis assesses how the resulting predictor rankings behave across partitions.

The posterior mean, 95 percent credible interval, posterior mean hazard ratio and interval are given, by feature. The likelihood that the coefficient is positive or negative is given as well as a direction. In the case of standardised numeric predictors, a one unit change results in the training set increase by 1 standard deviation. In the case of categorical variables, the effect is in relation to the reference category.



These posterior summaries describe conditional predictor associations and their uncertainty. Contribution rankings and SHAP attributions provide separate measures of explanatory importance \cite{ref40,ref53}.

The posterior probability is not considered to be frequentist statistical significance. A large posterior probability that a coefficient is positive indicates support for a positive hazard association under the given model and prior. It is not a standard p-value.

Posterior explainability is a complement of model agnostic explainability. Recent survival-XAI studies have shown that techniques such as LIME and SHAP can be used to explain individual survival predictions. They provide post-hoc attributions that complement the posterior summaries \cite{ref28,ref30,ref60}.

Based on this, the present thesis considers the Bayesian posterior per se to be the main explainability mechanism of the Bayesian models, and the main interpretable outputs to be the posterior hazard ratios and credible intervals.

\begin{table}[htbp]
\centering
\begin{tabular}{@{}p{0.33\linewidth}p{0.63\linewidth}@{}}
\toprule
Posterior quantity & Interpretation\\
\midrule
Coefficient distribution & Conditional log-hazard association\\
Draw-wise hazard ratio & $\exp(\beta_k)$ for each retained draw\\
95\% credible interval & Posterior uncertainty in the association\\
$\Pr(\beta_k>0\mid D)$ & Posterior support for a positive association\\
Contribution ranking & Relative magnitude of predictor contributions\\
\bottomrule
\end{tabular}
\caption{Bayesian posterior explainability outputs and their interpretation}
\label{_Toc245465019}
\end{table}

\section{Prior-Sensitivity Methodology}\label{prior-sensitivity-methodology}

To comprehend whether the conclusions made out of the posterior were sensitive to the prior family in the Bayesian survival model, a sensitivity analysis was carried out. Prior sensitivity is a key consideration in Bayesian analysis. When the observed data have little data or the parameters are not well defined the assumptions made in the prior may have an impact on the posterior inferences. Prior specification should balance the information encoded by the prior with the information contributed by the observed data \cite{ref40,ref56}.

Three prior families were compared in this study:

\begin{enumerate}
\def\labelenumi{\arabic{enumi}.}
\item
  Normal prior, which provides a common regularisation scale for the coefficients;
\item
  Horseshoe prior, with global-local shrinkage of possibly weak or noisy predictors; and
\item
  Spike-and-Slab prior, which explicates a mixture of concentrated spike and more diffuse slab.
\end{enumerate}

The sensitivity analysis compared both prediction and posterior behaviour under controlled changes to the modelling conditions. Training-information conditions used nested event/time-stratified subsets containing 25\%, 50\%, 75\% and 100\% of outer-training observations. Subsets were shared across priors, and a condition required at least 20 training events. Administrative censoring truncated follow-up at the outer-training 75th or 50th percentile, at full training size. These conditions used outer folds 1 and 2 in each repeat, giving ten intended paired outcomes per cohort and condition.

Design sensitivities used outer fold 1 in each repeat. They examined engineered GBSG2 predictors, published and engineered WHAS500 predictors, and METABRIC NPI-component, stage and treatment-augmented representations. The METABRIC stage condition used five training-only iterative posterior-sampling imputations; subject-level survival predictions were averaged across imputations before recalculating metrics. Baseline-hazard sensitivities used 4, 6 and 10 intervals. Prior-scale conditions halved or doubled the primary regularisation scales, with other factors held fixed. The full sensitivity design comprised 825 model--condition--fold outcomes and 885 underlying Bayesian fits, including the multiple-imputation fits.

Relative IBS changes were summarised as differences from the matching primary condition and as difference-in-differences between priors. Uncertainty used 2,000 paired three-stage bootstrap replicates, resampling repeats, folds within repeats and validation subjects within event strata. Coverage was reported alongside each estimate. These analyses assessed sensitivity to the information available and the specified modelling choices.

\subsection{Interpretability and SHAP Comparisons}\label{interpretability-shap-boundary-and-permutation-importance}

Predictive evaluation was kept separate from interpretability. SHAP is a popular post-hoc explanation framework that has been increasingly used for survival prediction, including to compare between machine learning survival models in cancer. Further studies have shown that model agnostic explanation methods can give individualised explanations for survival curves and time-dependent survival predictions \cite{ref28,ref30,ref60}.

The completed SHAP analysis targeted the linear log-risk score. For posterior-mean Bayesian coefficient $\overline\beta_j$, subject value $x_{ij}$ and training-background mean $\overline x_j$, exact interventional linear SHAP is
\begin{equation}
\phi_{ij}=\overline\beta_j(x_{ij}-\overline x_j).
\end{equation}
The explainer output was checked against this expression \cite{lundberg2017,shaplinear}. Same-model SHAP explained the Bayesian posterior-mean risk score; CoxPH SHAP supplied a secondary cross-model comparison.

Posterior contribution importance for source-variable group $g$ was
\begin{equation}
I_g=\frac{1}{N}\sum_i E\left[\left|\sum_{j\in g}x_{ij}\beta_j\right|\mid D\right].
\end{equation}
For both contribution and SHAP rankings, signed encoded terms were summed within a subject and clinical variable before taking absolute magnitude and averaging. This kept categorical predictors grouped by their clinical meaning.

Three sampling seeds and nested validation/background budgets of 25/50, 50/100 and 100/200 were used, capped at the available fold sizes. Selections were paired across priors. Successive budgets within a seed and adjacent seeds at the largest budget were compared using a Spearman threshold of 0.95 and at least four shared groups in the top five. At the largest budget, 2,000 hierarchical bootstrap replicates resampled repeats, folds and explanation seeds. Stable positive agreement required all evaluable stability comparisons to pass and a Spearman interval wholly above zero. Cohort/prior claims also required complete fold coverage. This separated posterior direction, cross-method rank agreement and sampling stability.

\subsection{Reproducibility and Computational Integrity.}\label{reproducibility-and-computational-integrity.}

Reproducibility was an important part of the computational analysis. Deterministic random seeds were derived from base seed 260817. All parameters were specified for the models, priors, MCMC settings, diagnostics and what was being evaluated. The Bayesian models had 2000 warmup iterations and 2000 posterior samples were drawn. The target acceptance probabilities were specified per model. For the Normal and Spike-and-Slab models the target acceptance probabilities were 0.95. For the Horseshoe model the target acceptance probabilities were 0.99. The Normal, Horseshoe, and Spike-and-Slab model settings were recorded explicitly to support reproducibility. This makes the calculations and comparisons reproducible. Reproducible Bayesian workflows focus on specifying the computational details, diagnostics and re-validating model assumptions \cite{ref39,ref40,ref53}.

The diagnostic summaries, fitted settings, preprocessing state, validation predictions and posterior summaries were retained for each outcome. Acceptance was determined by zero divergences, zero maximum-tree-depth hits, minimum BFMI at least 0.30, bulk ESS above 400, R-hat below 1.01 and finite diagnostic aggregates. Tail ESS and Monte Carlo standard error were recorded as supplementary summaries. Scientific reproduction checks compared all 450 original model--fold outcomes with restored computations, while allowing machine-dependent timing and archive hashes to differ.

\section{Supplementary Analyses}\label{supplementary-analyses}

Two supplementary analyses were specified separately and labelled post hoc. A sampler sensitivity repeated all 75 horseshoe outer fits with target acceptance 0.999, keeping the folds, likelihood, predictors, priors and diagnostic criteria fixed. These results were assessed separately from the original prior comparisons. A mathematical Holm bound examined the attainable adjusted values for the observed primary tests while allowing unavailable test values to range over their admissible values. This assessed the effect of incomplete test-family coverage without estimating missing model effects.

\section{Chapter Summary}\label{cross-check-of-chapter-and-methodology.}

This chapter describes the methods used to address the research questions. Bayesian survival models were evaluated across three clinical cohorts alongside statistical and machine-learning comparators. Three prior families were examined, and posterior uncertainty was assessed using MCMC diagnostics and Bayesian explainability summaries \cite{ref40}.

The methodology combines three clinical cohorts, three Bayesian prior families and three comparator model families within five repeats of five-fold validation. Posterior interpretation, explanation stability and sensitivity analyses address complementary aspects of the fitted models. The research questions are linked to these components below.

\begin{longtable}[]{@{}
  >{\raggedright\arraybackslash}p{(\columnwidth - 2\tabcolsep) * \real{0.2003}}
  >{\raggedright\arraybackslash}p{(\columnwidth - 2\tabcolsep) * \real{0.7997}}@{}}
\caption{Mapping of the research questions to the methodological components}\label{_Toc245465020}\\
\toprule\noalign{}
\begin{minipage}[b]{\linewidth}\raggedright
Research Question
\end{minipage} & \begin{minipage}[b]{\linewidth}\raggedright
Methodological Component
\end{minipage} \\
\midrule\noalign{}
\endhead
\bottomrule\noalign{}
\endlastfoot
RQ1 & Bayesian models versus CoxPH, RSF, and DeepSurv predictive evaluation \\
RQ2 & Normal, Horseshoe, and Spike-and-Slab posterior comparison \\
RQ3 & Posterior coefficients, hazard ratios, credible intervals, and directional probabilities \\
RQ4 & Diagnostic assessment and sensitivity to training information, censoring, predictors, baseline intervals and prior scales \\
\end{longtable}
